\section{Model III: Cost, Feasibility, and Minimax Regret}
\subsection{Model Description}
 
Model III turns the physical and comfort results into a decision. It checks whether each option has a valid model solution and whether it meets any active cost, electricity, or water limits. It then compares the options in three hot-weather scenarios. It selects the option with the smallest worst-case regret across the three scenarios.
 
\subsection{Variables}
\begin{table}[H]
  \centering
  \renewcommand{\arraystretch}{1.3}
  \begin{tabularx}{\textwidth}{p{0.20\textwidth}lY}
    \toprule
    \textbf{Variable} & \textbf{Symbol} & \textbf{Description} \\
    \midrule
    Equipment option & $i,j$ & Candidate option and index used to compare feasible alternatives. \\
    \rowcolor{gray!15} Weather scenario & $s$ & Peak-temperature, hot-humid, or hot-dry scenario. \\
    Hot-discomfort loss & $L_{i,s}$ & Modeled discomfort for option $i$ in scenario $s$ ($^\circ$C$\cdot$h). \\
    \rowcolor{gray!15} Scenario regret & $R_{i,s}$ & Excess loss over the best feasible option in that scenario ($^\circ$C$\cdot$h). \\
    Maximum regret & $R_{i,\max}$ & Largest scenario regret for an option ($^\circ$C$\cdot$h). \\
    \rowcolor{gray!15} Recommended option & $i^*$ & Feasible option that minimizes maximum regret. \\
    Equipment capital cost & $C$ & Upfront equipment cost in dollars. \\
    \bottomrule
  \end{tabularx}
\end{table}

\subsection{Primary Equations}
 
We use the minimax-regret decision criterion, applied here to discomfort loss~\cite{savage1951}. First define the set of physically valid options in weather scenario $s$ as $V_s$. Let $C_{i,s}$, $E_{i,s}$, and $m_{water,i,s}$ be the active cost, electricity use, and water use. Given limits on these resources, the scenario's feasible set is
\begin{equation}
 F_s=\left\{i\in V_s:
   C_{i,s}\le C_{cap},\quad
   E_{i,s}\le E_{cap},\quad
   m_{water,i,s}\le m_{cap}\right\}.
\end{equation}
We compare each resource use with the corresponding limit. The set of devices that can be recommended across all modeled weather is $F=\bigcap_s F_s$. Each scenario still has its own benchmark, $L_s^*=\min_{j\in F_s}L_{j,s}$. After observing the weather an operator could choose any device in $F_s$, so $L_s^*$ is the best attainable loss in that scenario. Subtracting that benchmark from the loss for a fixed option $i$ gives the regret
\begin{equation}
R_{i,s} = L_{i,s} - \min_{j\ \text{feasible}} L_{j,s}
\end{equation}
An option has zero regret if it is best in that scenario. Its worst regret is
\begin{equation}
R_{i,\max} = \max_{s} R_{i,s}
\end{equation}
The minimax recommendation is the robust option with the smallest upper bound on its scenario regrets. For a fixed option, the smallest valid bound $z$ is its largest scenario regret:
\begin{equation}
 \min_z\{z: z\ge L_{i,s}-L_s^* \text{ for every }s\}
 = \max_s(L_{i,s}-L_s^*)=R_{i,\max}.
\end{equation}
Minimizing that bound over the robust set gives
\begin{equation}
i^{*} \in \arg\min_{i\in F}\; R_{i,\max}
\end{equation}
We use equipment purchase prices as the cost measure for the comparative analysis.
 
\subsection{Choice of Decision Criterion}
 
TOPSIS ranks choices by their distance from an ideal point and a worst point. Its usual weighted formulation normalizes the criteria and assigns criterion weights~\cite{hwang_yoon1981}. We considered four criteria: cost, energy, water, and comfort. Those choices can strongly affect the result, especially when the units are different. A TOPSIS score can look precise even when the weights are subjective.
 
We use minimax regret to compare each choice with the best choice available under each weather scenario. The model screens options against resource limits and ranks the feasible options by discomfort regret. Each limit can stay in its original units.
 
At Grove Hall, the Large Evaporative Cooler leads on comfort and requires \$4,939 upfront, 2.50 kWh per hour, and up to about 84.1 kg of water in a modeled hour. Under illustrative limits of \$4,500, 1.50 kWh per hour, and 40 kg of water per hour, maximum regret among feasible choices is zero for the Fan, 0.49 SET degree-hours for the Small Evaporative Cooler, and 1.71 for no cooling.
 
\subsection[Key Assumptions]{Key Assumptions and Justifications}
 
\begin{table}[H]
  \centering
  \renewcommand{\arraystretch}{1.3}
  \begin{tabularx}{\textwidth}{lYp{0.66\textwidth}}
    \toprule
    \textbf{\#} & \textbf{Assumption} & \multicolumn{1}{c}{\textbf{Justification}} \\
    \midrule
    \nextassumption\label{assumption:model3-1} & The three selected weather scenarios represent the uncertainty of interest. & Peak, humid-hot, and dry-hot hours provide contrasting stresses. Each case tests an operating condition. \\
    \rowcolor{gray!15} \nextassumption\label{assumption:model3-2} & Regret is based on hot-discomfort loss. & We compare discomfort after screening options for resource feasibility. \\
    \nextassumption\label{assumption:model3-3} & An option must be feasible in every scenario to be robust. & This avoids recommending a device that fails under one supplied condition. \\
    \rowcolor{gray!15} \nextassumption\label{assumption:model3-4} & Purchase cost is the selected cost basis. & It is available and easy to compare. It omits maintenance and long-term utility costs. \\
    \nextassumption\label{assumption:model3-5} & Resource limits describe the venue. & Purchase cost, electricity, and water can each have a separate limit. \\
    \rowcolor{gray!15} \nextassumption\label{assumption:model3-6} & Scenarios enter the decision unweighted. & The largest regret across the three hours determines the ranking. \\
    \bottomrule
  \end{tabularx}
\end{table}